% Force directions, the three probe points A, B, C and the trajectory of the
% worked example.  Data: ch15-field.dat, ch15-example-traj.dat.
\begin{tikzpicture}
\begin{axis}[
  width=0.6\textwidth,axis equal image,
  xmin=-0.4,xmax=9.4,ymin=0.1,ymax=8.0,
  xlabel={$x$},ylabel={$y$},font=\footnotesize,
  grid=major,grid style={black!10},
]
\addplot[sbGray,thin,quiver={u=\thisrow{u},v=\thisrow{v},scale arrows=0.34,
  every arrow/.append style={-{Stealth[length=2.6pt,width=2pt]}}}]
  table[x=x,y=y]{figures/data/ch15-field.dat};
\addplot[fill=sbObstacle,draw=black!40,domain=0:360,samples=61,variable=\t] ({4+cos(\t)},{4+sin(\t)});
\addplot[dashed,black!50,domain=0:360,samples=73,variable=\t] ({4+3*cos(\t)},{4+3*sin(\t)});
\addplot[sbPath,line width=1.5pt] table[x=x,y=y]{figures/data/ch15-example-traj.dat};
\node[sbstart] at (axis cs:0,4.5) {};
\node[sbgoal] at (axis cs:8,4) {};
\node[sbannot,above] at (axis cs:8,4.15) {goal};
\node[sbannot,above] at (axis cs:6.1,6.55) {$\rho=\rho_0$};
\addplot[only marks,mark=*,mark size=1.8pt,sbRed] coordinates {(0.5,4) (2.5,4) (3,5.3)};
\node[sbannot,below,text=sbRed] at (axis cs:0.5,3.85) {A};
\node[sbannot,below,text=sbRed] at (axis cs:2.5,3.85) {B};
\node[sbannot,left,text=sbRed] at (axis cs:2.85,5.35) {C};
\end{axis}
\end{tikzpicture}
